II) \emph{Reducción al caso en que $B$ es una $A$-álgebra de tipo
finito que es un dominio entero y contiene a $A$.} --- Sea
$(B_\alpha)$ la familia filtrante creciente de las sub-$A$-álgebras
de tipo finito de $B$, y sea $C_\alpha$ la clausura
\oldpage[IV-2]{175}
integral de $A$ en $B_\alpha$. De la definición de clausura integral se sigue de
inmediato que $C$ es la unión de los $C_\alpha$. Si se pone
$B'_\alpha=B_\alpha\otimes_A A'$ y
$C'_\alpha=C_\alpha\otimes_A A'$, la planitud implica que
$C'_\alpha$ está contenido en $B'_\alpha$ y, por la misma razón,
$B'$ es la unión de la familia filtrante creciente de los
$B'_\alpha$, mientras que $C'$ es la unión de la familia filtrante
creciente de los $C'_\alpha$. Si se demuestra que $C'_\alpha$ es la
clausura integral de $A'$ en $B'_\alpha$, se seguirá de inmediato
que $C'$ es la clausura integral de $A'$ en $B'$.
